\section{引言：为什么需要 Scaling Laws}

假设你有一位非常有钱的朋友，给了你 100,000 张 H100 GPU 用一个月，目标是构建最好的开源大语言模型。前面的课程已经教会了你分布式训练框架、预训练数据准备、模型架构选择等基本要素。但面对如此庞大的资源，一个核心问题浮现了：\textbf{如何在不浪费算力的前提下，做出最优的工程决策？}

\begin{figure}[H]
\centering
\includegraphics[width=0.95\textwidth]{slides-images/slide-001.jpg}
\caption{课程大纲：Scaling Laws 的核心动机与学习目标\protect\footnotemark}
\end{figure}
\footnotetext{Slides 第2页。}

简单地``照搬别人的做法''（如 follow LLaMA 的配置）虽然安全，但无法真正推动前沿。如果你在一家前沿实验室，想要创新和优化，就需要一套\textbf{系统性的方法}来预测不同决策在大规模下的表现。

\begin{importantbox}{Scaling Laws 的核心价值}
Scaling Laws 的核心思想是：\textbf{通过训练一系列小模型来预测大模型的行为}。不再需要花费巨额算力直接调优大模型，而是在小规模上学习规律，然后外推到大规模，``一次搞定''大模型训练。
\end{importantbox}

传统深度学习的做法是：直接训练大模型、调超参数、发现效果不好、再调——这极其昂贵。新的``乐观主义''方法是：
\begin{enumerate}
    \item 在小规模上训练大量模型
    \item 从小模型中学习行为规律
    \item 外推到大规模，一次性训练出最优大模型
\end{enumerate}

\begin{figure}[H]
\centering
\includegraphics[width=0.95\textwidth]{slides-images/slide-002.jpg}
\caption{从小模型学习规律，外推到大模型——这是 Scaling Laws 的基本方法论\protect\footnotemark}
\end{figure}
\footnotetext{Slides 第3页。}

\subsection{本章小结}

Scaling Laws 为大语言模型的工程决策提供了一套基于数据驱动的科学方法。其核心前提是：模型在小规模上展示的趋势可以可靠地外推到大规模。这一思想贯穿本节课的所有内容。

%% ========== Part 2: 历史与背景 ==========

